\textbf{结论名称：} 全称量词命题 $\forall x\in M, p(x)$ 的否定

\textbf{适用条件：} 命题中含有全称量词 $\forall$，且谓词 $p(x)$ 可判定真值。

\textbf{变量说明：} $M$ 为非空集合（论域），$p(x)$ 为关于 $x$ 的命题函数（谓词），$\neg p(x)$ 表示 $p(x)$ 的否定。

\textbf{核心公式：}
\[
\boxed{\neg \big( \forall x \in M, \ p(x) \big) \ \iff \ \exists x \in M, \ \neg p(x)}
\]

\textbf{使用场景：} 否定全称命题，用于反证、逆否命题构造、逻辑化简，以及日常语言中“不是所有…都…”的数学化表达。

\textbf{简单例子：} 设 $M = \{1,2,3\}$，$p(x)$ 为 “$x > 2$”。原命题 $\forall x\in M, x>2$ 为假（因为 $1\le 2$）。其否定：$\exists x\in M, \neg(x>2)$，即 $\exists x\in M, x\le 2$，显然为真（$x=1$ 或 $2$ 均为证据），验证等价关系成立。